\chapter{Introduction}
\label{chap:intro}

Cooperative vehicular applications depend on the timely wireless dissemination of safety-relevant information among nearby vehicles. \gls{v2x} communication enables vehicles to exchange messages directly with other vehicles, infrastructure, and pedestrians, and \gls{ieee}~802.11 remains an important reference technology for this purpose because it supports direct, decentralized exchanges without continuous infrastructure support \cite{IEEE_80211p_Survey}. In particular, IEEE~802.11p adapts the standard for vehicular operation, typically at 5.9~GHz with 10~MHz channelization, and relies on fully distributed channel access: there is no central scheduler that can reserve transmission opportunities for urgent traffic.

Among the messages exchanged in such networks, crash alerts occupy a special position. They are rare, but when they occur, their usefulness depends almost entirely on timely delivery to nearby vehicles. Contention-based channel access, however, exposes this urgent crash-related traffic to the same medium competition faced by routine transmissions: as offered load rises, delay and jitter become unstable, backoff freezing can postpone emergency frames while the channel is busy, and \gls{qos} markings are not handled consistently across deployments \cite{RFC8325_WiFi_QoS, Continuous_Backoff_Freezing_Li}. A crash alert that arrives late competes for the channel exactly when the surrounding traffic is most likely to be dense, which is also when its delivery matters most.

Within IEEE~802.11, two standard mechanisms are central to traffic differentiation. At the network layer, \gls{dscp} marking expresses traffic criticality; at the link layer, \gls{edca} translates that criticality into differentiated medium access through access categories with distinct contention parameters \cite{IEEE80211}. The resulting wireless behavior, however, depends on the local policy that maps DSCP values into 802.11 user priorities and then into EDCA queues, and default mappings are ambiguous in practice \cite{RFC8325_WiFi_QoS}. Moreover, EDCA provides statistical preference rather than guarantees: higher-priority categories obtain access earlier on average, but all categories still contend randomly, and urgent traffic improves only at the expense of lower-priority flows.

Recent literature pursues stronger contention control through slot-based schedulers inspired by \gls{tsn}, preemptive or cooperative \gls{mac} schemes, and cellular sidelink resource allocation \cite{TSNCtl_Feraudo, PRP_MAC_Li, CFC_MAC_Linn, CV2X_Sidelink_Allocation}. These approaches confirm the value of tighter contention control, but they rely on synchronization, infrastructure, or richer protocol machinery, assumptions that are difficult to satisfy in a decentralized vehicular ad hoc network. This leaves a practical gap for a local, event-driven design that isolates one crash-triggered high-priority flow against ordinary background traffic and quantifies both the protection obtained and the penalty imposed. Staying within standard IEEE~802.11p mechanisms keeps such a design deployable: no clock synchronization, infrastructure, or protocol changes are required, and adoption can proceed per node.

This dissertation addresses that gap with a deliberately narrow study: a focused crash-aware prioritization strategy built from standard mechanisms rather than a new MAC protocol. All vehicles exchange periodic \gls{be} multicast traffic, while one designated crash vehicle emits temporary alert bursts marked with DSCP~46 and mapped to the \gls{vo} access category. On top of this two-class model, five MAC policies are compared: a plain \gls{dcf} baseline, EDCA with explicit DSCP-to-Voice mapping, and three adaptive extensions (\texttt{stable}, \texttt{guarded}, and \texttt{emergency}) that temporarily suppress or drop competing BE traffic while alerts are active. The adaptive profiles are fixed design points that map a protection-versus-cost frontier, aiming at better statistical latency under congestion rather than deterministic guarantees.

\section{Research Question and Scope}
\label{sec:intro:question}

The central question investigated in this dissertation is:

\begin{quote}
	Can crash-triggered traffic obtain better wireless service than ordinary traffic in Veins/\gls{inet} simulations, and what cost does that protection impose on ordinary traffic?
\end{quote}

The scope is intentionally minimal. The study considers a single crash source, a two-class BE/VO traffic distinction, and one-hop multicast dissemination on a highway corridor, with no access points, roadside units, or centralized coordination. It does not attempt to reproduce a full \gls{its} stack, standardized cooperative-awareness message schedules, or multi-class traffic orchestration. All findings are obtained by simulation in \gls{omnet}/Veins/\gls{inet}, and all observed prioritization is statistical contention shaping: no deterministic slotting, synchronization, or provable worst-case bounds are claimed.

\section{Objectives}
\label{sec:intro:objectives}

The general objective of this dissertation is to evaluate, through simulation, whether crash-triggered Voice traffic can obtain measurably better wireless service than ordinary Best Effort traffic in decentralized vehicular Wi-Fi, using only local, standard-compliant MAC mechanisms, and to quantify the cost that this protection imposes on ordinary traffic.

The specific objectives are:

\begin{enumerate}
	\item To implement a crash-triggered two-class traffic model in Veins/\gls{inet}, in which all vehicles generate periodic BE multicast traffic and a designated crash vehicle emits DSCP~46 alert bursts mapped to the VO access category through an explicit classifier.
	\item To design and implement a lightweight adaptive MAC controller that extends EDCA with a local three-mode protection mechanism, exposed as a ladder of five policies with increasing aggressiveness, from plain DCF to emergency preemption.
	\item To evaluate the five policies across a matrix of vehicle densities (10 and 100 vehicles) and offered-load profiles, using class-specific delay, jitter, multicast reach, and MAC-loss metrics.
	\item To characterize the resulting protection-versus-cost frontier and derive engineering guidance on when adaptive suppression should be enabled.
\end{enumerate}

\section{Contributions}
\label{sec:intro:contributions}

This dissertation makes the following primary contributions:

\begin{enumerate}
	\item A lightweight adaptive MAC controller built solely from standard mechanisms---explicit DSCP~46-to-Voice mapping plus a local three-mode extension of EDCA---applied to a crash-triggered two-class traffic model in which ordinary vehicular multicast traffic remains Best Effort while a designated crash node injects temporary Voice alert bursts.
	\item An open-source Veins/\gls{inet} simulation artifact \cite{aguiar2026_veins_qos} comparing a plain IEEE~802.11p DCF baseline, an EDCA baseline with explicit mapping, and three adaptive profiles (\texttt{stable}, \texttt{guarded}, \texttt{emergency}) across density and load matrices.
	\item Engineering guidance on \emph{when} adaptive suppression should be enabled: under heavy congestion, emergency preemption compresses crash-alert tail delay beyond explicit EDCA (approximately $-49\%$ in the 95th-percentile Voice delay) at the cost of much higher BE delay and explicit BE losses, whereas at light contention bounded suppression mainly penalizes BE---supporting event-triggered escalation rather than always-on suppression.
\end{enumerate}

\section{Document Structure}
\label{sec:intro:structure}

The remainder of this dissertation is structured as follows. Chapter~2 provides the necessary background on IEEE~802.11p, contention-based channel access, EDCA, and the DSCP-to-Voice mapping path. Chapter~3 surveys related work on vehicular QoS and stronger contention-control alternatives. Chapter~4 formalizes the system model: the highway scenario, the two-class traffic abstraction, the crash event timeline, and the five MAC policies. Chapter~5 details the implementation in \gls{omnet}/Veins/\gls{inet} and presents the evaluation results and their discussion. Finally, Chapter~6 outlines the work plan and concludes.
